\documentclass[10pt,a4paper]{amsart}
\usepackage[margin=19mm]{geometry}
\usepackage{amsmath,amssymb,mathtools}
\usepackage[scr=rsfs,cal=euler]{mathalfa}
\usepackage{xfrac}
\usepackage[unicode,hidelinks,bookmarks=false]{hyperref}
\newcommand{\ie}{i.e.\ }
\newcommand{\cf}{cf.\ }
\newcommand{\resp}{respectively\ }
\newcommand{\Subsection}{\subsection}
\providecommand{\nobreakdash}{\nobreak\hbox{-}}
\setlength{\emergencystretch}{2em}
\multlinegap0pt
\usepackage{bm}
\usepackage{bbm}
\usepackage{dsfont}
\usepackage{xspace}
\usepackage{cancel}
\usepackage{mathtools}
\usepackage{amsthm}
\makeatletter
\@ifundefined{theorem}{\newtheorem{theorem}{Theorem}}{}
\@ifundefined{lemma}{\newtheorem{lemma}[theorem]{Lemma}}{}
\@ifundefined{proposition}{\newtheorem{proposition}[theorem]{Proposition}}{}
\makeatother
\title[Singularity formation and self-shrinkers for the parabolic B]{Singularity formation and self-shrinkers for the parabolic Bernoulli problem}
\author{Hunter Liu, Sebastian Munoz, Zihui Zhao}
\date{Source: arXiv:2608.22574v1 (CC BY 4.0)}
\begin{document}
\maketitle

\noindent\emph{Source and licence.} Singularity formation and self-shrinkers for the parabolic Bernoulli problem. Version arXiv:2608.22574v1 (CC BY 4.0), distributed under the Creative Commons Attribution 4.0 International licence, as recorded in the arXiv licence metadata for this exact version and shown on the abstract page https://arxiv.org/abs/2608.22574v1. Full rights notice: licenses/arxiv-2608.22574-CC-BY-4.0.txt.\par\smallskip

\noindent\textit{Editorial scope.} Counted unit: the Proposition whose source label is \texttt{prop:annulus-width} in the article (source0.tex lines 2235--2241, source label \texttt{prop:annulus-width}), delivered together with the article's own complete original proof of it (source0.tex lines 2243--2404, one \texttt{proof} environment of 162 lines). No mathematics is written, repaired or re-derived: statement and proof are the article's own text line for line. Required context retained uncounted: eq:profile (lines 88--93); eq:ode (lines 144--146); prop:odesol (lines 951--962); eq:annular-scaled-notation (lines 1926--1934); eq:annular-y-equation (lines 1936--1939); eq:hermite-kappa (lines 1971--1975); lem:kappa-bounds (lines 1990--1995); lem:annular-phase (lines 2045--2053); lem:annular-barrier (lines 2130--2137). Every definition, display and label of those lines is retained unchanged. Omitted from this package and not redistributed: 1--87, 94--143, 147--950, 963--1925, 1935--1935, 1940--1970, 1976--1989, 1996--2044, 2054--2129, 2138--2234, 2242--2242, 2405--2968. Nothing inside a retained excerpt is omitted or altered: every retained body is delivered exactly as the article's lines are written, so no omission receipt of any kind accompanies this package. Citations: the retained excerpts carry no citation command at all, so no bibliography entry of the article is transcribed and no reference is redistributed. Reference adaptation: none was needed and none was made; every reference inside the delivered unit resolves inside it, so no unresolved reference or citation is produced. Printed slips kept literal: no printed word, symbol, hypothesis or number of the counted statement or of its proof was altered. Layout and class: the article's own class is \texttt{article} with packages including amsmath, amssymb, amsthm, mathrsfs, geometry, enumitem, xcolor, soul, authblk, hyperref; the wrapper is the lane's pinned \texttt{amsart} editorial wrapper at 10pt on A4, which restates the theorem-like environments the retained text needs and adds the article's own macro definitions that the retained text needs (none), with extra wrapper packages (bm, bbm, dsfont, xspace, cancel, mathtools). The delivered numbering is the unit's own. Assets: no retained span contains a figure environment, a graphics-inclusion command, a PDF-page inclusion, a table or a TikZ picture; no image file is copied into this package, the package declares no asset dependency and no image file is redistributed with it. Licence: version arXiv:2608.22574v1 of this article is distributed under the Creative Commons Attribution 4.0 International licence, as recorded in the arXiv licence metadata for this exact version and shown on the abstract page https://arxiv.org/abs/2608.22574v1. A full rights notice is delivered at licenses/arxiv-2608.22574-CC-BY-4.0.txt. Characters: every retained excerpt is pure ASCII, so the delivered engine, XeLaTeX with its default Latin Modern text font, reports no missing character.
\begin{equation}\label{eq:profile}
\left\{\begin{array}{ll}
	\Delta\phi - \tfrac{1}{2}\,y \cdot \nabla\phi + \tfrac{1}{2}\,\phi = 0 & \text{in } \Omega:= \{\phi>0\}, \\
	|\nabla\phi| = 1 & \text{on } \partial\Omega.
\end{array}   \right.\tag{P}
\end{equation}

\begin{equation}\label{eq:ode}
    \phi''(r) + \left(\frac{n-1}{r} - \frac{r}{2}\right) \phi'(r) + \frac{1}{2}\,\phi(r) = 0.\tag{RP}
\end{equation}

\begin{proposition}\label{prop:odesol}
If $n>1$, there exist radii $0 < r_- < r_+$ and a function $\phi \in C^2((r_-,r_+)) \cap C^1([r_-,r_+])$ such that:
\begin{enumerate}
\item[\textup{(i)}] $\phi > 0$ on $(r_-,r_+)$ and $\phi(r_-) = \phi(r_+) = 0$;
\item[\textup{(ii)}] $\phi$ solves the radial profile ODE \eqref{eq:ode}, i.e.
\[
\phi''(r) + \left(\frac{n-1}{r} - \frac{r}{2}\right)\phi'(r) + \frac{1}{2}\,\phi(r) = 0, \qquad r \in (r_-,r_+);
\]
\item[\textup{(iii)}] the Bernoulli conditions hold: $\phi'(r_-) = 1$ and $\phi'(r_+) = -1$.
\end{enumerate}
The function defined by $\phi(y) := \phi(|y|)$ for $y$ in the annulus $\{r_- < |y| < r_+\}$ and $\phi \equiv 0$ otherwise is the unique radial solution of~\eqref{eq:profile} compactly supported on the annulus.
\end{proposition}

\begin{equation}\label{eq:annular-scaled-notation}
    \mu:=\frac{n-1}{2},
    \qquad
    a:=\frac{r_-}{\sqrt2},
    \qquad
    b:=\frac{r_+}{\sqrt2},
    \qquad
    y(t):=\frac1{\sqrt2}\phi(\sqrt2\,t),
\end{equation}

\begin{equation}\label{eq:annular-y-equation}
    y''-2h(t)y'+y=0
    \qquad\text{on }(a,b),\tag{D}
\end{equation}

\begin{equation}\label{eq:hermite-kappa}
    Y''-2xY'+Y=0,
    \qquad
    Y(0)=1,\quad Y'(0)=0;\tag{H}
\end{equation}

\begin{lemma}\label{lem:kappa-bounds}
The constant defined by \eqref{eq:hermite-kappa} satisfies
\begin{equation}\label{eq:kappa-bounds}
    \frac43<\kappa^2<\frac{23}{17}.
\end{equation}
\end{lemma}

\begin{lemma}\label{lem:annular-phase}
The scaled profile $y$ in \eqref{eq:annular-y-equation} has a unique critical point $c\in(a,b)$. Moreover,
\[
    c\geq\sqrt{2\mu}
    \qquad\text{and}\qquad
    ab<2\mu;
\]
in particular, $r_-r_+<2(n-1)$.
\end{lemma}

\begin{lemma}\label{lem:annular-barrier}
We have
\[
    a>\sqrt{2\mu}-\kappa
    \qquad\text{and}\qquad
    b>\sqrt{2\mu}+\kappa.
\]
\end{lemma}

\begin{proposition}\label{prop:annulus-width}
Let $\phi$ be the annular profile from Proposition~\ref{prop:odesol}, with support $[r_-,r_+]$. For every $n\geq2$,
\[
    r_+-r_-<\sqrt{14};
\]
equivalently, in the scaled variables \eqref{eq:annular-scaled-notation}, $b-a<\sqrt7$.
\end{proposition}

\begin{proof}
Write $p:=ab$ and
\[
    \varrho(t):=t^{2\mu}e^{-t^2/2},
\]
so that $\varrho'=-2h\varrho$ and the equation \eqref{eq:annular-y-equation} becomes divergence form $(\varrho\,y')'+\varrho\,y=0$. By Lemma~\ref{lem:annular-phase}, $p<2\mu$ for every $n\geq2$; for $n\geq3$, Lemmas \ref{lem:annular-barrier} and \ref{lem:kappa-bounds} complement this with a lower bound, so that $p$ lies in the window
\begin{equation}\label{eq:product-window}
    2\mu-\frac{23}{17}
    <2\mu-\kappa^2
    =\bigl(\sqrt{2\mu}-\kappa\bigr)\bigl(\sqrt{2\mu}+\kappa\bigr)
    <p<2\mu.
\end{equation}

Suppose for contradiction that $b-a\geq\sqrt7$. For fixed product $p$, the endpoints of the interval with product $p$ and width $d$ are $\frac12\bigl(\sqrt{d^2+4p}\mp d\bigr)$, the first decreasing and the second increasing in $d$. Hence the interval $[\alpha,\beta]$ determined by
\[
    \beta-\alpha=\sqrt7,
    \qquad
    \alpha\beta=p,
\]
satisfies $[\alpha,\beta]\subset[a,b]$, and $\alpha>0$.

\emph{Step 1: reduction to a trial function.} Let $\lambda_1$ be the first eigenvalue of the Dirichlet problem
\begin{equation}\label{eq:width-eigen}
    -(\varrho\,\eta')'=\lambda\,\varrho\,\eta
    \quad\text{on }(\alpha,\beta),
    \qquad
    \eta(\alpha)=\eta(\beta)=0.
\end{equation}
Since $y>0$ on $(a,b)$ and solves $-(\varrho y')'=\varrho y$ with $y(a)=y(b)=0$, the first Dirichlet eigenvalue on $(a,b)$ is $1$. Hence domain monotonicity and $[\alpha,\beta]\subset[a,b]$ give $\lambda_1\geq1$. It therefore suffices to show $\lambda_1<1$.

The Liouville substitution $v=\varrho^{1/2}\eta$ transforms \eqref{eq:width-eigen} into the Dirichlet problem for $-v''+V_n\,v=\lambda v$ on $(\alpha,\beta)$, with
\begin{equation}\label{eq:width-potential}
    V_n(t)=\frac{t^2}{4}+\frac{\mu(\mu-1)}{t^2}-\mu-\frac12,
\end{equation}
so it suffices to exhibit $\psi\in H^1_0((\alpha,\beta))$ with
\begin{equation}\label{eq:width-rayleigh}
    \frac{\displaystyle\int_\alpha^\beta\bigl[(\psi')^2+V_n\,\psi^2\bigr]\,dt}
         {\displaystyle\int_\alpha^\beta\psi^2\,dt}<1 .
\end{equation}
Write
\[
    L:=\sqrt7,
    \qquad
    m:=\frac{\alpha+\beta}2=\sqrt{p+\frac74},
    \qquad
    t=m+x,
    \quad
    -\frac L2\leq x\leq\frac L2 .
\]

\emph{Step 2: the case $n\geq3$.} Take
\[
    \psi(x):=\cos\frac{\pi x}{L},
\]
and let $\langle F\rangle:=\int F\psi^2\,dx\big/\!\int\psi^2\,dx$ denote the normalized averages. Elementary integration gives
\begin{equation}\label{eq:width-cos-moments}
    \frac{\int(\psi')^2}{\int\psi^2}=\frac{\pi^2}{7},
    \qquad
    \langle t^2\rangle=m^2+\vartheta\,\frac{L^2}4=p+C,
    \qquad
    C:=\frac73-\frac7{2\pi^2},
    \qquad
    \vartheta:=\frac13-\frac2{\pi^2}.
\end{equation}
Since $|x|\leq\frac L2<m$, pairing $x$ with $-x$ gives the convergent expansion
\[
    \langle t^{-2}\rangle
    =\frac1{m^2}\biggl\{1+\sum_{j\geq1}(2j+1)\Bigl\langle\Bigl(\frac xm\Bigr)^{2j}\Bigr\rangle\biggr\},
\]
and $\langle x^{2j}\rangle\leq\bigl(\tfrac L2\bigr)^{2j-2}\langle x^2\rangle=\vartheta\bigl(\tfrac L2\bigr)^{2j}$ yields, after summing the resulting series,
\begin{equation}\label{eq:width-inverse-upper}
    \langle t^{-2}\rangle\leq J_\vartheta(p),
    \qquad
    J_\vartheta(p):=\frac4{4p+7}\left(1+\vartheta\,\frac{7(6p+7)}{8p^2}\right)
    =\frac{4(1-\vartheta)}{4p+7}+\frac\vartheta p+\frac{7\vartheta}{2p^2}.
\end{equation}
Since $\mu(\mu-1)\geq0$ for $n\geq3$, the Rayleigh quotient in \eqref{eq:width-rayleigh} is at most
\begin{equation}\label{eq:width-Q}
    \mathcal Q(p):=\frac{\pi^2}7+\frac{p+C}4+\mu(\mu-1)\,J_\vartheta(p)-\mu-\frac12 .
\end{equation}

We claim that $\mathcal Q$ is increasing throughout the window \eqref{eq:product-window}. The second form of $J_\vartheta$ in \eqref{eq:width-inverse-upper} shows $J_\vartheta''>0$, and $\pi^2<10$ gives $\vartheta<\frac2{15}$. Since $\partial_\vartheta J_\vartheta'(p)=\frac{16}{(4p+7)^2}-\frac1{p^2}-\frac7{p^3}<0$,
\[
    4\,\mathcal Q'(p)
    =1+4\mu(\mu-1)\,J_\vartheta'(p)
    \geq 1+(n-1)(n-3)\,J_{2/15}'(p),
\]
with strict inequality when $n\geq 4$, and the right-hand side is increasing in $p$. For $n=3$ it equals $1$. For $n\geq4$, set $B:=n-4\geq0$; clearing denominators at the left endpoint $p=2\mu-\frac{23}{17}$ of the window gives
\[
    1+(n-1)(n-3)\,J_{2/15}'\!\left(2\mu-\frac{23}{17}\right)
    =\frac{P(B)}{15\,(17B+28)^3\,(68B+231)^2}>0,
\]
where
\[
    P(B):=475067448\,B^4+2882697837\,B^3+5070969134\,B^2+2143250690\,B+10362030 .
\]
Hence $\mathcal Q(p)<\mathcal Q(2\mu)$ throughout the window, by the upper bound in \eqref{eq:product-window}. Finally, since $\partial_\vartheta J_\vartheta=\frac1p+\frac7{2p^2}-\frac4{4p+7}>0$, the bounds $\vartheta<\frac2{15}$ and $\pi^2<10$ give
\[
    4\bigl(1-\mathcal Q(2\mu)\bigr)
    >n+5-\frac{3233}{420}-(n-1)(n-3)\,J_{2/15}(n-1)
    =\frac{2272\,n^2-8555\,n+9027}{420\,(n-1)(4n+3)}>0,
\]
where the quadratic in the numerator has negative discriminant. Thus the Rayleigh quotient of $\psi$ is at most $\mathcal Q(p)<\mathcal Q(2\mu)<1$, and Step~1 applies.

\emph{Step 3: the case $n=2$.} Here $\mu=\frac12$, and two features of Step~2 change. First, the coefficient $\mu(\mu-1)=-\frac14$ of the inverse-square term in \eqref{eq:width-potential} is now negative, so an upper bound such as \eqref{eq:width-inverse-upper} no longer controls the Rayleigh quotient, and we need a lower bound on the average of $t^{-2}$ instead. Second, the lower bound in \eqref{eq:product-window} becomes vacuous, since $2\mu-\kappa^2<0$; the only property of the product $p$ we use is $0<p<2\mu=1$. Put $k:=\pi/L$ and take the two-mode trial function
\[
    \psi(x):=\cos(kx)-\frac1{10}\sin(2kx),
\]
which vanishes at $x=\pm\frac L2$. Let
\[
    K:=\frac{\int(\psi')^2}{\int\psi^2},
    \qquad
    T:=\frac{\int t^2\psi^2}{\int\psi^2},
    \qquad
    S:=\frac{\int t^{-2}\psi^2}{\int\psi^2}.
\]
Orthogonality and direct integration give
\begin{equation}\label{eq:width-two-mode-KT}
    K=\frac{104}{101}\cdot\frac{\pi^2}7,
    \qquad
    T=m^2+\frac7{12}-\frac{2807}{808\,\pi^2}-\frac{640\,mL}{909\,\pi^2}.
\end{equation}
Since $\pi^2>9$ and $L<3$ give $\frac{640L}{909\pi^2}<1<2m$, the expression for $T$ is increasing in $m$; as $m^2=p+\frac74\leq\frac{11}4$, we may evaluate it at $m^2=\frac{11}4$, where $mL=\frac{\sqrt{77}}2$. Using $\pi<\frac{22}7$ and $\sqrt{77}>\frac{877}{100}$,
\[
    K=\frac{104}{707}\,\pi^2<\frac{104}{707}\cdot\frac{484}{49}=\frac{50336}{34643}<\frac{1453}{1000},
    \qquad
    T<\frac{11}4+\frac7{12}-\frac{49}{484}\left(\frac{2807}{808}+\frac{14032}{4545}\right)<\frac{267}{100}.
\]
For a lower bound on $S$, write
\[
    E(x):=\cos^2(kx)+\frac1{100}\sin^2(2kx),
\]
so that $\psi^2=E-\frac15\cos(kx)\sin(2kx)$ and $\int\psi^2=\int E$. The omitted cross term is odd; on $\bigl(0,\frac L2\bigr)$ it is nonpositive, while $(m+x)^{-2}-(m-x)^{-2}<0$ there, so the cross term contributes nonnegatively to $\int t^{-2}\psi^2$. Pairing $x$ with $-x$ as in Step~2 and keeping the first three (positive) terms of the expansion therefore gives
\begin{equation}\label{eq:width-two-mode-S}
    S\geq\frac1{m^2}+\frac{3M_2}{m^4}+\frac{5M_4}{m^6},
    \qquad
    M_2:=\frac{\int x^2E}{\int E},
    \qquad
    M_4:=\frac{\int x^4E}{\int E},
\end{equation}
where direct integration gives
\[
    M_2=\frac7{12}-\frac{2807}{808\,\pi^2},
    \qquad
    M_4=49\left(\frac1{80}-\frac{401}{1616\,\pi^2}+\frac{4803}{3232\,\pi^4}\right).
\]
Both $M_2$ and $M_4$ are positive, so the right-hand side of \eqref{eq:width-two-mode-S} is decreasing in $m^2$ and may be evaluated at $m^2=\frac{11}4$, where it equals
\[
    \frac4{11}+\frac{48}{121}\,M_2+\frac{320}{1331}\,M_4 .
\]
As a function of $y:=\pi^2$, this expression has derivative $\frac{14\,(41303\,y-336210)}{134431\,y^3}$, which is positive for $y\geq9$; substituting the lower bound $\pi>\frac{333}{106}$ therefore yields
\[
    S>\frac{267896812791500}{551004452874117}>\frac{243}{500}.
\]
For $n=2$ the Rayleigh quotient in \eqref{eq:width-rayleigh} is
\[
    K+\frac T4-\frac S4-1
    <\frac{1453}{1000}+\frac{267}{400}-\frac{243}{2000}-1
    =\frac{999}{1000}<1,
\]
and Step~1 applies.
\end{proof}

\end{document}
